\section{Finish the checksum and release the descriptor}
A successful EOF is the only path that returns a finalized checksum. Complement the
running state once, close the descriptor with syscall 16, and return status zero.
Closing uses \texttt{v0} to select a service, so retain the checksum somewhere else
until the close has completed. A checksum held only in \texttt{v0} would be replaced
by the close service number before return.

A read error follows a different result path but has the same resource obligation:
close the descriptor, return checksum zero, and return status two. An unsuccessful
open returns status one without closing an invalid result. The close service supplies
no separate success result in this educational ABI. Do not interpret the service
selector left in a register as a returned close status.

The file routine calls an update routine, and the table method also calls the table
initializer. Each call replaces \texttt{ra}. Save the incoming return address in the
frame before making either call, then restore it before returning to the caller.
Likewise retain the buffer pointer, capacity, descriptor, method and running state
across calls rather than trusting caller-saved temporaries to survive.

The illustrated frame occupies 64 bytes and saves seven full saved registers plus
\texttt{ra}. Eight-byte stores and loads retain their high halves. The original
stack alignment is maintained because the frame size is a multiple of eight. A
routine that uses a different set of registers may use a different frame, provided
it meets the same preservation and return contract.

\subsection*{Exercise: Audit all exits}
Draw the control-flow paths for unsuccessful open, immediate EOF, several successful
reads followed by EOF, and several reads followed by an error. Mark where a valid
descriptor first exists, where it is closed, and where the checksum becomes finalized.
For each path, identify the state that must survive a nested update call.
